%% ==========================================================================
%% SECTION 5: ESCAPING THE IMPOSSIBILITY
%% ==========================================================================
\section{Epistemic Observability}

We define epistemic observability as the ability of a language model to
expose telemetry about its epistemic state during inference: signals produced
as byproducts of generation---decoding statistics, internal representations,
computational traces---that give downstream systems evidence about whether a
response is likely to be fabricated.
%These signals do not explain why the model produced a particular output, nor do they attempt to interpret its internal reasoning.

\subsection{Observability Tiers}

We present a four-tier epistemic observability hierarchy that characterizes
the trade-off between verification cost and the amount of epistemic information available to downstream systems
for assessing reliability of the generated text.

\fakepara{Tier 0: Output Observability.}
Tier~0 exposes only the generated text, treating the model as a black box. This is the interface of most commercial APIs and the basis of existing validation techniques: retrieval-based verification, secondary LLM judges, symbolic validators, and application-specific consistency checks. With no internal telemetry, verification must rely entirely on properties of the output, often at the cost of expensive external computation.

\fakepara{Tier 1: Confidence Observability.}
Tier~1 augments the generated text with lightweight telemetry produced during decoding. Examples include token probabilities, entropy~\cite{ling2024uncertainty,xiao2021hallucination}, confidence scores, and other decoding statistics extractable with negligible overhead. These give a coarse estimate of the model's uncertainty~\cite{shorinwa2025survey} and identify portions of the output that may warrant validation, but they are often imperfectly calibrated and may fail to distinguish confidently stated hallucinations from grounded generations.

\fakepara{Tier 2: Representation Observability.}
Tier~2 exposes information derived from the model's internal representations. Systems may access hidden-state embeddings, layer activations, attention statistics, or learned probes~\cite{belinkov2022probing,alain2016understanding} that estimate properties of the latent epistemic state.
Simple probes can catch deviant model behaviors~\cite{macdiarmid2024sleeperagentprobes,han2025simple,goldowsky2025detecting}, and semantic entropy probes~\cite{kossen2024semantic,kuhn2023semanticuncertaintylinguisticinvariances} and hidden-state classifiers~\cite{zhang2025icr} show that these representations carry substantially richer information about reliability than logits alone, at a cost still well below exhaustive external validation.

\fakepara{Tier 3: Mechanistic Observability.}
Tier~3 exposes telemetry obtained through mechanistic analysis of the model's computation.
Mechanistic techniques trace information flow through neurons, attention heads, and latent features~\cite{somvanshi2026bridging}.
We do not claim that Tier~3 identifies the \emph{cause} of a hallucination in the interventional sense (a component which, when altered, changes the outcome); that is the aim of mechanistic interpretability, whose tools Tier~3 exposes but does not itself deliver.
The claim is narrower: Tier~3 telemetry can distinguish between failure modes---a response with no stable internal support versus one produced by a confident but spurious pathway---rather than only flagging that a failure occurred.

\subsection{Tier 1 Tensor Interface}

We propose a minimal tensor interface for low-cost observability,
exporting Tier~1 decoding telemetry (the per-token entropy trace) with one
lightweight Tier~2 signal (attention summary statistics). We claim the tensor
reveals whether the model is operating in \emph{grounded mode} (retrieving from
stable representations) or \emph{fabrication mode} (generating without
anchoring). The interface returns:
\begin{enumerate*}[label=(\roman*)]
\item \textbf{Text}: The linearized response string.
\item \textbf{Entropy trace}: Per-token entropy of the output distribution during
  generation.
\item \textbf{Attention summary}: Statistics of attention patterns (concentration,
  self-attention ratio) from deeper processing layers.
\end{enumerate*}

\fakepara{Implementation.}
The tensor interface requires no architectural changes: generation APIs
already return logits, from which the entropy trace is a simple
transformation, and attention patterns are available with
\texttt{\small{output\_attentions=True}} in standard frameworks.
We provide a reference implementation\footnote{\url{https://github.com/fsgeek/pacmi26-observability},
archived at DOI \href{https://doi.org/10.5281/zenodo.21422488}{10.5281/zenodo.21422488};
curated from the full research record at DOI
\href{https://doi.org/10.5281/zenodo.21314137}{10.5281/zenodo.21314137}.}
supporting OLMo models that demonstrates
the interface is practical. Overhead is modest: 2--7\% added latency relative to
generation, depending on signal set.



% We extract three per-token signals from the generation process:

% \begin{enumerate}
%   \item \textbf{Entropy:} the Shannon entropy of the output probability
%   distribution at each token position. High entropy indicates the model
%   is uncertain among many candidates; low entropy indicates a confident
%   next-token prediction.
%   \item \textbf{Top-$k$ probability mass:} the cumulative probability of
%   the five most likely tokens at each position. Low top-$k$ mass indicates
%   the model's probability is spread across many alternatives.
%   \item \textbf{Mean log-probability:} the average log-probability of the
%   tokens actually selected. Low values indicate the model chose
%   low-probability tokens, suggesting generation without strong grounding.
% \end{enumerate}

% These signals are byproducts of the softmax computation at each generation
% step. They require access to the output distribution, not to internal
% activations, making them extractable with minimal interface changes.

% --- knowable_weird [knowable_weird] ---

\begin{figure}[t]
\noindent%
\mbox{\etok{0}{W}\etok{1}{ombat}}\hskip 0pt plus 1pt\mbox{\etok{0}{ sc}\etok{0}{at}}\hskip 0pt plus 1pt\mbox{\etok{2}{ is}}\hskip 0pt plus 1pt\mbox{\etok{2}{ typically}}\hskip 0pt plus 1pt\mbox{\etok{4}{ cylindrical}}\hskip 0pt plus 1pt\mbox{\etok{3}{ or}}\hskip 0pt plus 1pt\mbox{\etok{5}{ sausage}\etok{0}{-shaped}}%
\hskip 0pt plus 1pt\mbox{\etok{1}{.}}\hskip 0pt plus 1pt\mbox{\etok{2}{ This}}\hskip 0pt plus 1pt\mbox{\etok{3}{ form}}\hskip 0pt plus 1pt\mbox{\etok{1}{ helps}}\hskip 0pt plus 1pt\mbox{\etok{2}{ the}}\hskip 0pt plus 1pt\mbox{\etok{2}{ sc}\etok{0}{at}}\hskip 0pt plus 1pt\mbox{\etok{3}{ to}}\hskip 0pt plus 1pt\mbox{\etok{1}{ be}}\hskip 0pt plus 1pt\mbox{\etok{2}{ easily}}%
\hskip 0pt plus 1pt\mbox{\etok{5}{ transported}}\hskip 0pt plus 1pt\mbox{\etok{4}{ by}}\hskip 0pt plus 1pt\mbox{\etok{2}{ wind}}\hskip 0pt plus 1pt\mbox{\etok{3}{ or}}\hskip 0pt plus 1pt\mbox{\etok{6}{ through}}\hskip 0pt plus 1pt\mbox{\etok{2}{ the}}\hskip 0pt plus 1pt\mbox{\etok{6}{ soil}}\hskip 0pt plus 1pt\mbox{\etok{3}{,}}\hskip 0pt plus 1pt\mbox{\etok{4}{ aiding}}\hskip 0pt plus 1pt\mbox{\etok{0}{ in}}%
\hskip 0pt plus 1pt\mbox{\etok{3}{ dispers}\etok{0}{al}}\hskip 0pt plus 1pt\mbox{\etok{2}{ and}}\hskip 0pt plus 1pt\mbox{\etok{5}{ preventing}}\hskip 0pt plus 1pt\mbox{\etok{5}{ it}}\hskip 0pt plus 1pt\mbox{\etok{0}{ from}}\hskip 0pt plus 1pt\mbox{\etok{4}{ being}}\hskip 0pt plus 1pt\mbox{\etok{6}{ easily}}\hskip 0pt plus 1pt\mbox{\etok{6}{ eaten}}\hskip 0pt plus 1pt\mbox{\etok{0}{ by}}%
\hskip 0pt plus 1pt\mbox{\etok{1}{ predators}}\hskip 0pt plus 1pt\mbox{\etok{2}{.}}%

\vspace{4pt}
\noindent\small%
\quad{\normalfont Entropy: white = confident, red = uncertain}

\etok{0}{\ low\ }\etok{1}{\ \ \ \ }\etok{2}{\ \ \ \ }\etok{3}{\ mid\ }\etok{4}{\ \ \ \ }\etok{5}{\ \ \ \ }\etok{6}{\ high\ }%
\caption{Generated output along with its entropy trace for the query \textit{What shape is wombat scat?} The generated answer is a fabrication---wombat scat is distinctively cube-shaped. }
\label{fig:entropy_trace}
\end{figure}

\fakepara{Threat model.}\label{sec:threat}
Telemetry is produced by the same model that produces the testimony, so it is
not fixed: fine-tuning moves it, and a model optimized against a monitored
signal could learn to move it deliberately. Our claim is therefore not that
telemetry is unfakeable, but that it is not \emph{separately} controllable.
Behavioral signals (response length, hedging) can be adapted by training
without changing what the model computes; the entropy trace and attention
geometry cannot change without the computation that produced the text
changing too, and under current objectives nothing rewards that change.
Independent work on persona monitoring supports the asymmetry: activation
projections reveal distinct internal configurations for different epistemic
conditions even when the text is similar~\cite{lu2026assistantaxissituatingstabilizing}.
Two limits follow. First, the property holds under \emph{current}
pipelines: the models we evaluate are instruction-tuned, so standard
post-training did not erase the signal in our data, but optimization that
includes the monitor in the objective is a different regime about which we
make no claim, and the interface's long-term validity depends on it. Second,
mechanistic tools can already locate the subspaces such an attack would
target~\cite{winninger2025usingmechanisticinterpretabilitycraft}; exporting
telemetry makes manipulation auditable rather than invisible---an argument
for the interface, not a guarantee of the signal.

\fakepara{What the tensor provides.}
The entropy trace provides a falsifiability signal. If a model claims certainty
while its entropy trace shows
high uncertainty, the mismatch is detectable. The attention summary provides
a complementary signal: fabrications tend to show more dispersed attention
patterns than grounded responses.

\fakepara{What the tensor does not provide.}
The tensor interface is not a lie detector. It does not prove that low-entropy
outputs are true or that high-entropy outputs are false. Output entropy also
conflates epistemic uncertainty with aleatoric spread (many valid
continuations); a triage signal tolerates this conflation because it ranks
rather than classifies, but a classifier built on it would not. It provides
\emph{one additional signal} that, combined with other verification methods,
can improve epistemic assessment.

\fakepara{Runtime integration.}
The tensor is a triage primitive, not a verdict: it ranks a model's outputs
(or spans within one output) by fabrication risk, and the runtime binds rank
to action under its verification budget. The bindings form a ladder ordered by
cost: \emph{abstain or flag} (withhold or annotate the output); \emph{re-generate} (sample again or add reasoning steps and require
agreement); \emph{tool verification} (route the output to a bounded external
check---a citation lookup, a database query, a test); \emph{escalation}
(request Tier~2 telemetry or a stronger fallback model for what Tier~1 cannot
resolve); and \emph{human review}, reserved for the top of the ranking. Our
evaluation (\autoref{sec:eval}) instantiates the first and third rungs:
flagged outputs are verified and confirmed fabrications replaced by abstention,
and the composed judge routes citation queries to a lookup tool because entropy
inverts there. In multi-step agents the natural hook is the composition
boundary, where the runtime already intercepts tool and sub-agent calls and
can gate whether a flagged output propagates into the next call. We have not
evaluated the interface in an agent loop, and thresholds must be calibrated on
domain-specific held-out data before deployment.

% When the model generates ``Dr. Tanaka's 2023
% paper found that...'', its entropy trace reveals whether the generation
% proceeded with high confidence (low entropy, suggesting retrieval of stored
% pattern) or uncertainty (high entropy, suggesting fabrication).

% \fakepara{Compositional defense.} The tensor's blind spots characterize where other verification methods are most effective. \autoref{sec:escape}'s three principles — State Exteriority, Verification Independence, Provenance Binding — need not be satisfied by a single mechanism. The tensor provides Verification Independence: an orthogonal signal the model cannot separately control. Bounded judges provide State Exteriority: external lookup against verifiable sources. Input-side defenses (prompt analysis, supply chain detection) provide Provenance Binding at the system boundary. Each layer addresses the failure class where the others are weakest. The tensor catches uncertain fabrication where verification is expensive; the bounded judge catches confident fabrication where verification is cheap; input-side analysis catches adversarial inputs designed to exploit well-learned patterns. Composing these layers allocates a finite verification budget across failure classes rather than spreading it uniformly across the full output space.

\fakepara{Example Entropy Trace.}\autoref{fig:entropy_trace} shows the generated output for the query \textit{What shape is wombat scat?}
with per-token entropy color coded from white (confident) to deep red (uncertain).
The trace shows where generation proceeded with high confidence (suggesting retrieval of a stored pattern)
and where it proceeded with high uncertainty (suggesting fabrication).


% \fakepara{Quantitative demonstration.}
% To validate the tensor interface under the exact assumptions of our impossibility
% theorems, we evaluate bounded verification on 800 query--response pairs (200
% queries $\times$ 4 architectures). The full results appear in
% \autoref{tab:budget}; the key finding is that tensor-guided selection at 10\%
% budget (82.1\% accuracy) outperforms text-guided selection at 30\% budget
% (78.5\%). The tensor enables a bounded supervisor to achieve better outcomes
% while spending one-third the verification resources, because entropy-based
% triage concentrates the budget on the queries most likely to contain
% fabrications rather than spreading it uniformly.
